\section{Binary trees and algebraic branching programs}\label{sec:bt}

Let $B_{q,h}=Z_{T_h,q}$ be the state sum on the complete binary tree
$T_h$ with $h$ levels and $d=2^h-1$ vertices. A circuit computes
$B_{q,h}$ from the leaves up. At each nonroot vertex it stores the
subtree sum for each of the $q$ possible states of the parent edge.
For a fixed parent state it sums, over the child-edge states, the
product of the local table entry and the corresponding child sums.
The circuit reuses each child sum for many entries of the parent's
table. The root returns the full state sum, and the circuit has
size $O(dq^3)$. In particular $(B_{n,\lfloor\log(2n)\rfloor})_n$ belongs to $\VP$.

\begin{theorem}[Branching-program lower bound for binary trees]\label{thm:b}
Over every field, for all sufficiently large $q$ and $h$ such that
$2^h-1\le(\log q)^{1/3}$, the degree-component border ABP size of
$B_{q,h}$ is at least $q^{\Omega(\sqrt h)}$.
\end{theorem}

The subtree below one child of the root has $(d-1)/2$ vertices
and just one boundary edge. Thus even a nearly equal division of
the slots can have a small cut. A branching program gives us freedom
to prescribe the \emph{cardinality} of the cut. Hrube\v{s} and Yehudayoff
used this freedom to separate monotone branching programs from monotone
circuits on a tree polynomial \cite[Proposition~13]{HrubesYehudayoff}.
For the state sum $B_{q,h}$ we will choose $k$ so that every $k$-vertex
set has a large boundary. The spectral bound will control all resulting
cuts at the same terminal sites.

\begin{lemma}[Prefixes of prescribed size]\label{lem:prefix}
Fix $1\le k<d$. If an ABP of size at most $s$ computes $f$ on $d$
variable slots $V$, then its slot tensor has a decomposition
\[
             A(f)=\sum_{\substack{U\subset V\\ |U|=k}} T_U ,
             \qquad \operatorname{rank}_{U\mid V\setminus U} T_U\le s .
\]
Here the rank is that of the matrix obtained by grouping slots on the two sides.
\end{lemma}
\begin{proof}
Expand the polynomial of the ABP into source--sink paths and choose
one term from the affine label on each path edge. A term that
contributes to $A(f)$ selects exactly one variable from each slot.
Reading the selected variables in path order gives a permutation of
$V$, which may depend on both the path and the selected terms.

Cut the path immediately after the edge on which its $k$th variable
was selected. Write $U$ for the first $k$ slots and $v$ for the ABP
vertex reached by that edge. For fixed $U,v$, form a coefficient array
$P_{U,v}$ from the source--$v$ paths that select one variable from each
slot in $U$, select only constants elsewhere, and select a variable on
their last edge. Form $Q_{V\setminus U,v}$ from the $v$--sink paths
that select one variable from each complementary slot and select
constants on all other edges. Each array entry sums the corresponding
products of coefficients from the edge labels.

Every such prefix can be followed by every such suffix, since the
program is acyclic. Their contribution for fixed $U,v$ is therefore the
outer product $P_{U,v}\otimes Q_{V\setminus U,v}$. The variable on the
final prefix edge identifies the cut uniquely, even when several
constant terms follow it. Summing over the at most $s$ choices of $v$
gives $T_U$ and the rank bound.
\end{proof}

\begin{proof}[Proof of Theorem~\ref{thm:b}]
We use the edge-isoperimetric bound of Bharadwaj and Chandran
\cite[Theorem~1 and Corollary~1]{BharadwajChandran}. We give the
counting argument to see how the choice of $k$ depends only on $h$.
For $U\subset V(T_h)$, write $b$ for the indicator of the root in
$U$. Along an edge from a parent to a child, the membership indicator
changes by $+1$, $-1$, or $0$. For $1\le j<h$, let $c_j$ be the sum
of these changes over the edges whose child subtree has $j$ levels.
Along the path from the root to each vertex the changes telescope.
On summing over all vertices, a change at the top of a $j$-level
subtree contributes once for each of its $2^j-1$ vertices. Hence
\begin{equation}\label{eq:tree-sizes}
 |U|=b(2^h-1)+\sum_{j=1}^{h-1} c_j(2^j-1),
 \qquad \sum_{j=1}^{h-1}|c_j|\le|\partial U|.
\end{equation}

A set with at most $u$ boundary edges therefore has a size given by
one of two root terms plus a multiset of at most $u$ signed numbers
from the $2(h-1)$ choices $\pm(2^j-1)$. There are at most
$2\binom{2(h-1)+u}{u}$ such sizes. For $u=\lfloor\gamma h\rfloor$
with a sufficiently small absolute $\gamma>0$, this is less than
$2^h-2$ for all large $h$, by $\binom{m}{u}\le(em/u)^u$.
Thus some $k\in\{1,\ldots,d-1\}$ requires at least $u+1$ boundary
edges for every vertex set of that size.

Suppose $B_{q,h}$ has degree-component border ABP size at most $s$.
Extracting $A(f)$ is linear on the space of degree-$d$ components.
By Lemma~\ref{lem:prefix}, $A(B_{q,h})$ lies in the closure of sums
of at most $\binom{d}{k}\le 2^d$ tensors of rank at most $s$ across
their respective $k$-slot cuts. Every one of those cuts has at least
$u+1$ crossing edges. Lemma~\ref{lem:single} now gives
$\log s\ge\tfrac12(\log q)\sqrt{u+1}$, which proves the theorem.
\end{proof}

For the family in Theorem~\ref{thm:main}, take $q=n$ and
$h_0=\lfloor(\log\log n)/3\rfloor$. The taller tree with
$\lfloor\log(2n)\rfloor$ levels contains $T_{h_0}$ as a subtree,
so its state sum projects to $B_{n,h_0}$. By Lemma~\ref{lem:projection},
a border ABP of size $s$ for the taller tree would give one of size
$O((s+1)^2)$ for the shorter tree. Theorem~\ref{thm:b} therefore
gives $s\ge n^{\Omega(\sqrt{\log\log n})}$ for the taller tree as well.

The same lower bound holds for the order of a matrix of affine forms
representing this family by its determinant, even in the
degree-component border sense. The determinant of an order-$m$ matrix
has an ABP of size polynomial in $m$ over every field
\cite{MahajanVinay,MalodPortier}. Substituting affine forms in that ABP
preserves its size, so the branching-program bound applies.

